\section{Step 20: Require exact command completion}
Reaching pixel 4096 proves that commands cover the image. It does not prove
that the payload ends there. Compare the byte cursor with the exclusive payload
end and reject any leftover byte, including an otherwise valid extra command.
This makes the record extent unambiguous for the file reader.

Fill \texttt{record\_exact}. For the worked delta, the second command skips
4095 pixels and changes none. It consumes only its four-byte header and brings
both cursors to their required ends. The checkpoint that follows is the first
place from which either application loop may run. A valid prefix followed by a
bad late command must never reach it.

\paragraph{Observe} Prepare \texttt{s20-1} and stop before
\texttt{record\_validation\_ready}. At \texttt{record\_validation\_ready}, the byte cursor is
source plus 27 and the pixel cursor is 4096. All frame and context bytes are
still unchanged.

\Needspace{7\baselineskip}
\paragraph{Check} What must hold before applying a delta?
\begin{itemize}
\item[A.] Only its first command must fit.
\item[B.] Both the pixel coverage and payload consumption must be exact.
\item[C.] The checksum field must be rewritten.
\end{itemize}
